\section{AgentQ: de la «captura de pantalla del fallo» a trayectorias de recuperación de las que se puede aprender}

Este tramo es la parte de toda la clase que el ritmo de las demo oculta con más facilidad. El ponente muestra primero fallos comunes y después tareas de reservación, calendario y vuelos. El verdadero punto didáctico no es que «el agent puede reservar en un restaurante», sino cómo el sistema busca rutas alternativas, corrige errores, mantiene varias decisiones a la vez y recupera las trayectorias exitosas y fallidas como datos de entrenamiento.

\subsection{Primero, admitir que el Agent no es perfecto}

El collage de fallos incluye validación de formularios, pagos, horarios no disponibles y páginas de error. Le recuerda al lector que la retroalimentación de un entorno real casi nunca es un reward limpio, sino mensajes de error, resultados vacíos, páginas a medio completar y restricciones ocultas. La motivación de investigación de AgentQ es justamente que el agent siga buscando en esos estados, en lugar de dar por terminada la tarea tras un solo fallo.

Esta subsección continúa la distribución de evaluación del capítulo anterior: REAL nos dice en qué sitios se concentran los fallos, y AgentQ pregunta cómo actuar después de fallar. Al leer la figura de apertura, conviene separar primero los errores en «recuperables» y «no recuperables»: un campo faltante en un formulario suele poder corregirse, un pago rechazado puede requerir la intervención del usuario, y una identidad o un destinatario equivocados pueden obligar a terminar la tarea. Solo si se define antes el límite de la recuperación, el algoritmo de búsqueda no confundirá «seguir intentando» con el comportamiento correcto.

\lecturefigure{slide-19-agentq-failure-collage.jpg}{AgentQ abre mostrando tipos de fallo en páginas web reales}{00:17:59}

\lecturefigure{slide-20-agentq-title.jpg}{AgentQ: un marco de búsqueda y aprendizaje para Agent autónomos}{00:18:08}

\begin{importantbox}{El fallo debe estructurarse}
Un fallo se divide, como mínimo, en observation error, planning error, tool/action error, environment constraint y policy/safety rejection. Registrar solo «la tarea falló» no permite decidir si hay que cambiar el prompt, las herramientas, los datos, la recompensa o los permisos.
\end{importantbox}

\subsection{Tarea del restaurante: búsqueda, retroceso y conclusión}

El giro clave de la demo de reservación es que la hora objetivo no está disponible. El Agent consulta primero OpenTable, luego usa Google para obtener información complementaria y finalmente regresa al sitio original para completar la reservación. Los cuatro estados siguientes no están para mostrar una animación, sino para explicar un proceso de recuperación: la acción falla, se busca entre herramientas, se reubica y se completa con confirmación.

\lecturefigure{slide-21-reservation-prompt.jpg}{Objetivo del usuario: restaurante, hora y número de personas especificados}{00:18:11}

\lecturefigure{slide-22-agent-web-actions.jpg}{El Agent ejecuta acciones de reservación en la página web}{00:18:15}

\lecturefigure{slide-23-google-fallback.jpg}{Tras el fallo en el sitio original, recurre a Google para obtener información de disponibilidad}{00:18:19}

\lecturefigure{slide-24-adaptable-site-formats.jpg}{Regresa entre sitios y se adapta a distintas estructuras de página}{00:18:22}

\lecturefigure{slide-25-reservation-confirmed.jpg}{Tarea completada: la confirmación de la reservación se vuelve una poscondición verificable}{00:18:25}

\begin{knowledgebox}{Lectura de la figura: cómo validar la cadena de recuperación}
Primero, confirma si el agent reconoce la restricción del entorno «las 10 PM no están disponibles»; luego, revisa si la búsqueda entre sitios introdujo confusiones con restaurantes homónimos, fechas o zonas horarias; por último, verifica que la página de confirmación coincida con la intención original. La tarea solo cuenta como exitosa si la página final, la orden estructurada y el objetivo del usuario están alineados.
\end{knowledgebox}

\subsection{Self-correction y decisiones paralelas}

El siguiente grupo de figuras desglosa la capacidad en tres niveles: poder corregirse a sí mismo, poder manejar varias decisiones a la vez y poder reconocer una hora de calendario poco razonable y reprogramarla. Aquí, «self-correction» no debe entenderse como que el modelo diga «me equivoqué», sino como una máquina de estados que detecta que falló una poscondición, genera una acción alternativa y conserva las restricciones de la tarea original.

\lecturefigure{slide-26-self-correcting-behavior.jpg}{Árbol de comportamiento del Agent: tras una rama fallida, vuelve a elegir ruta}{00:18:28}

\lecturefigure{slide-27-multiple-decisions.jpg}{Mantener varias decisiones paralelas dentro de una misma tarea}{00:18:31}

\lecturefigure{slide-28-calendar-prompt.jpg}{Tarea de calendario: programar una reunión de una hora mañana a las 3}{00:18:32}

\lecturefigure{slide-29-calendar-rescheduled.jpg}{Reconoce que las 3 AM no son razonables y reprograma a las 3 PM}{00:18:39}

\begin{warningbox}{La «corrección automática» es lo que con más facilidad excede los permisos}
Si el usuario escribe «3», ¿cambiar de 3 AM a 3 PM es una inferencia razonable o una modificación no autorizada? En una tarea de calendario de bajo riesgo puede resolverse con una confirmación; en una tarea de alto riesgo, la corrección propuesta debe presentarse al usuario, no ejecutarse en silencio.
\end{warningbox}

\subsection{Skills routing y contexto de acciones}

«right skills at the right time» enfatiza el enrutamiento de herramientas. La tarea de vuelos requiere interpretar ciudades, fechas, viaje sencillo, preferencias de precio y de asiento, y después elegir la herramienta de búsqueda. El enrutador no es solo un clasificador: también debe decidir qué restricciones entran en los parámetros de la herramienta, cuáles se reservan para el filtrado posterior y si debe cambiar de herramienta tras un fallo.

Los casos anteriores del restaurante y del calendario muestran sobre todo la recuperación; esta sección pasa a la \textbf{selección de capacidades}. Que un agent tenga diez herramientas no significa que deban exponerse todas al modelo al mismo tiempo; cuantas más herramientas hay, más fácil es que ocurran nombres parecidos, parámetros en conflicto y usos indebidos de permisos. Al leer las figuras, conviene fijarse en la fidelidad con la que las restricciones se conservan a lo largo de la cadena: si el usuario dice «preferably one-way» y «window seat», la etapa de búsqueda puede relajar temporalmente el precio, pero no puede olvidar la preferencia de asiento en la recomendación final, y mucho menos cambiar en silencio la fecha de salida porque alguna herramienta respondió más rápido.

\lecturefigure{slide-30-right-skills-right-time.jpg}{Invocar la skill correcta en el momento correcto es un problema de orquestación del Agent}{00:18:41}

\lecturefigure{slide-31-flight-prompt.jpg}{La solicitud de vuelo incluye fecha, viaje sencillo, precio y preferencia de asiento}{00:18:42}

\lecturefigure{slide-32-flight-window-seat.jpg}{En los resultados de búsqueda sigue filtrando asientos de ventanilla en clase económica}{00:18:48}

\lecturefigure{slide-33-booking-success-improvement.jpg}{Mejora en el booking task success rate presentada en clase}{00:18:54}

Esta serie de demo también expone costos ocultos más allá de la tasa de éxito final. La búsqueda entre sitios aumenta la carga de páginas y los token, la corrección del calendario requiere explicar la ambigüedad horaria y el filtrado de vuelos mantiene varias restricciones blandas a la vez; una solicitud aparentemente sencilla de «hazlo por mí» incluye en realidad seguimiento de estado, cambio de herramientas, conservación de preferencias y verificación de resultados. Si el sistema solo optimiza la etiqueta de éxito, puede aprender a hacer clic de forma más agresiva o a ignorar las restricciones blandas. Los datos de entrenamiento deben incluir, como parte del resultado, la fidelidad completa a la intención del usuario.

\teachervoice{00:17:58--00:20:47, el ponente dice explícitamente que AgentQ no es un sistema perfecto; estas demo sirven para mostrar fallos, retrocesos, autocorrección, decisiones múltiples y skill routing. Las notas advierten: los resultados de la clase corresponden a tareas y configuraciones de modelo específicas, y no deben generalizarse como confiabilidad en cualquier sitio web.}

\begin{importantbox}{Una trace mínima reproducible de Agent}
\begin{lstlisting}[caption={Campos mínimos que debe guardar una trayectoria de Agent}]
trace = {
  "goal": user_goal,
  "initial_state": state_snapshot,
  "steps": [
    {"observation": obs, "thought_summary": plan,
     "action": action, "tool_result": result,
     "postcondition": check, "cost": cost}
  ],
  "outcome": outcome,
  "user_confirmations": confirmations,
  "rollback": rollback_record
}
\end{lstlisting}
\end{importantbox}

\subsection{Resumen de la sección}

Las demo de AgentQ no son publicidad de «reservación automática», sino un conjunto de casos de recuperación ante fallos: el entorno impone restricciones, el agent busca rutas alternativas, la máquina de estados verifica las poscondiciones, el enrutador elige la skill y la trayectoria se guarda como material de aprendizaje. El siguiente capítulo formaliza estos fenómenos como MCTS, process feedback y preference learning.
